\documentclass[10pt,a4paper]{amsart}
\usepackage[margin=19mm]{geometry}
\usepackage{amsmath,amssymb,mathtools}
\usepackage[scr=rsfs,cal=euler]{mathalfa}
\usepackage{xfrac}
\usepackage[unicode,hidelinks,bookmarks=false]{hyperref}
\newcommand{\ie}{i.e.\ }
\newcommand{\cf}{cf.\ }
\newcommand{\resp}{respectively\ }
\newcommand{\Subsection}{\subsection}
\providecommand{\nobreakdash}{\nobreak\hbox{-}}
\setlength{\emergencystretch}{2em}
\multlinegap0pt
\theoremstyle{plain}
\newtheorem{theorem}{Theorem}
\newtheorem{lemma}[theorem]{Lemma}
\theoremstyle{definition}
\newtheorem{definition}{Definition}
\newtheorem{ass}{Assumption}
\def \E {\mathbb{E}}
\newcommand{\Norm}[1]{\left\|#1\right\|}
\def \a {\mathbf{a}}
\def \e {\mathbf{e}}
\def \g {\mathbf{g}}
\def \q {\mathbf{q}}
\def \u {\mathbf{u}}
\def \v {\mathbf{v}}
\def \w {\mathbf{w}}
\def \x {\mathbf{x}}
\def \y {\mathbf{y}}
\def\eqref#1{equation~\ref{#1}}
\setcounter{secnumdepth}{0}
\title[Better Convergence Guarantees for Sign-Based Momentum Methods]{Better Convergence Guarantees for Sign-Based Momentum Methods: the MVSM-v2 last-iterate convergence guarantee (printed as Theorem 4.4)}
\author{Wei Jiang, Dingzhi Yu, Sifan Yang, Wenhao Yang, Zechao Li, Lijun Zhang}
\date{Source: arXiv:2507.12091v2 (CC BY 4.0)}
\begin{document}
\maketitle

\noindent\emph{Source and licence.} Better Convergence Guarantees for Sign-Based Momentum Methods. Version arXiv:2507.12091v2 (CC BY 4.0), distributed under the Creative Commons Attribution 4.0 International licence, http://creativecommons.org/licenses/by/4.0/, as recorded in the arXiv licence metadata for this exact version and shown on the saved abstract page https://arxiv.org/abs/2507.12091v2. Full licence notice: licenses/arxiv-2507.12091-CC-BY-4.0.txt.\par\smallskip

\noindent\textit{Editorial scope.} Counted unit: the last-iterate convergence guarantee for the distributed majority-vote signSGD-with-momentum method MVSM-v2, printed in the article as Theorem 4.4 (source0.tex lines 1003--1008, source label \texttt{thm4}), delivered together with the article's complete original proof of it (source0.tex lines 1324--1430, the article's own appendix section headed \emph{Proof of Theorem 4.4}). The article writes this proof as a detached appendix section rather than in a \texttt{proof} environment or as a run-in paragraph, so the proof block delivered here is exactly that section, heading line included, running to the last line before the next section. The proof is complete and finishes on the statement's own rate $O(d^{1/4}T^{-1/4})$; no line of it is omitted, shortened, substituted or re-flowed, and no line of the statement is either.

Required context retained uncounted: the article's bounded-below assumption with the initial suboptimality gap (lines 812--813, printed as Assumption 1); the article's per-node smoothness assumption (lines 909--913, printed as Assumption 6) and its per-node bounded-noise assumption (lines 914--920, printed as Assumption 7); the article's bounded-gradient assumption (lines 921--923, printed as Assumption 8); the article's definition of the unbiased sign operation $S_R$ with its display (lines 934--942, printed as Definition 4.1) and the remark recording the unbiasedness $E[S_R(\v)]=\v/R$ (line 943); the article's momentum estimator and its two-stage sign update (lines 945--952); and the article's definition of the MVSM-v2 server rule with its display (lines 998--1001). These four assumptions and that notation are exactly what the statement invokes and what the proof uses; nothing else from the article is needed for them.

Editorial formatting only, in addition: an AMS \texttt{amsart} preamble that restates the macro definitions the retained lines use -- the article's own $\E$, $\Norm{\cdot}$ and lower-case bold-vector macros $\a$, $\e$, $\g$, $\q$, $\u$, $\v$, $\w$, $\x$, $\y$ (copied from the article's own preamble), and the article's own definition of \texttt{\textbackslash eqref} -- together with the article's own theorem-environment declarations restated without the article's per-section counter and with an independent counter for the retained definition. Consequently the retained environments are numbered sequentially here: the four retained assumptions print as Assumptions 1, 2, 3 and 4, the retained definition prints as Definition 1, and the counted theorem prints as Theorem 1, whereas the article prints them as Assumptions 1, 6, 7 and 8, Definition 4.1 and Theorem 4.4. The wrapper prints no section number either (the article carries this proof as its unnumbered appendix section D), so the retained proof heading appears here as PROOF OF THEOREM 1. Every cross-reference inside the retained text resolves to this wrapper's numbering, which is the only numbering change; the equations of the retained proof are likewise numbered anew in order of appearance. No citation command and no \texttt{thebibliography} occur in any retained line, so this package carries no bibliography block and no bibliography entry. Outside the retained spans and therefore not delivered: the article's remark that follows the statement, comparing the rate with earlier work, and its short paragraph headed Core of the analysis; both carry citation commands, neither is part of the theorem environment or of the proof section, and the delivered proof re-derives the virtual update that the second of them only summarises, so neither carries a step the delivered proof uses. No figure, no image-inclusion line and no drawing code occur in any retained line, so no artwork is delivered or needed. The article's own style file \texttt{iclr2027\_conference.sty}, its \texttt{natbib} and \texttt{fancyhdr} style files and its \texttt{.bst} file are neither shipped in this package nor used to build it; the package ships the article's concatenated source verbatim as \texttt{sources/181671/source0.tex} and the retained excerpts are exact line ranges of it. No mathematical statement, hypothesis, estimate or conclusion has been rewritten, repaired or added, and printed slips, if any, are preserved literally.
\begin{ass}\label{ass:below} $f_{*}=\inf_x f(x) > -\infty$ and $f(\x_1) - f_{*} \leq \Delta_f$ for the initial solution $\x_1$.
\end{ass}

\begin{ass}(Smoothness on node $j$)\label{smooth-j} For each node $j\in[n]$, we suppose
    \begin{align*}
        \Norm{\nabla f_j(\x)-\nabla f_j(\y)}\le L\Norm{\x-\y}.
    \end{align*}
\end{ass}
\begin{ass}\label{ass:3++}  (Bounded noise on node $j$) For each node $j\in[n]$, we have
\begin{equation*}
\begin{split}
 \mathbb{E}_{\xi}\left[\left\|\nabla f_j(\x;\xi) -\nabla f_j(\mathbf{x})\right\|^{2}\right] \leq \sigma^{2}.
\end{split}
\end{equation*} 
\end{ass}

\begin{ass}(Bounded gradients)\label{bg1}
    For each node $j\in[n]$, we assume $\sup_{\x} \Norm{\nabla f_j(\x;\xi)}_\infty\le G$.
\end{ass}

\begin{definition}
    For any vector \( \v \) with \( \|\v\|_{\infty} \leq R \), the function $\operatorname{S_\textit{R}}(\v)$ is defined component-wise by:
    \begin{align}\label{mapping}
    [\operatorname{S_\textit{R}}(\v)]_k = \begin{cases}
        1, & \text{with probability } \frac{R+[\v]_k}{2R}, \\ \\
        -1, & \text{with probability } \frac{R-[\v]_k}{2R}.
    \end{cases}
    \end{align}
\end{definition}

\textbf{Remark:} This operation provides an unbiased estimate of \( {\v}/{R} \), since \( \E[\operatorname{S_\textit{R}}(\v)] = {\v}/{R} \). 

We can now introduce our majority vote signSGD with momentum updates. First, we use the momentum gradient estimator at each node $j$ as follows:
\begin{align}
    \v_t^j = (1-\beta) \v_{t-1}^j + \beta \nabla f_j(\x_t;\xi_t^{j}),
\end{align}
where $\beta$ is the momentum parameter. Next, by communicating the gradient estimators with the unbiased sign operation $\operatorname{S_\textit{G}}(\cdot)$, we update the decision variable as follows:
\begin{align}
    \x_{t+1} = \x_t - \eta \operatorname{Sign}\left(\frac{1}{n} \sum_{j=1}^{n} \operatorname{S_\textit{G}}(\v_t^j) \right).
\end{align} 

Although the above theorems achieve better convergence rates than previous methods, they do not converge to zero as $T$ increases. To address this issue, we replace the sign operation in the server with the unbiased sign operation $\operatorname{S_1}(\cdot)$ as defined in equation~(\ref{mapping}) with $R = 1$. The revised formulation for the update is:
\begin{align}\label{modi}
    \v_t = \operatorname{S_1}\left( \frac{1}{n}\sum_{j=1}^n \operatorname{S_\textit{G}}\left({\v}_t^j\right)\right).
\end{align}

\begin{theorem}\label{thm4}
Under Assumptions~\ref{ass:below}, \ref{smooth-j}, \ref{ass:3++}, and~\ref{bg1}, set $\beta\in(0,1]$ and $\eta=\mathcal{O}(1/\sqrt{dT})$. Then MVSM-v2 satisfies
\begin{align*}
 \E\|\nabla f(\x_\tau)\|_2=\mathcal O\left(\frac{d^{1/4}}{T^{1/4}}\right).
\end{align*}
\end{theorem}

\section{Proof of Theorem~\ref{thm4}}\label{app:mvsm-v2}
Define that $\bar\g_t =\frac1n\sum_{j=1}^n\nabla f_j(\x_t;\xi_t^j)$, $\bar\v_t=\frac1n\sum_{j=1}^n\v_t^j$, $\q_t^j=\operatorname{S_\textit{G}}(\v_t^j)$, $\a_t=\frac1n\sum_{j=1}^n\q_t^j$ and we can know that $\E[\bar\g_t]=\nabla f(\x_t)$. By Assumption~\ref{bg1} and induction, we have
\begin{align*}
 \|\v_1^j\|_\infty\le G,\qquad
 \|\v_t^j\|_\infty\le(1-\beta)\|\v_{t-1}^j\|_\infty
       +\beta\|\nabla f_j(\x_t;\xi_t^j)\|_\infty\le G.
\end{align*}
Every worker sign is therefore well defined. Also $\|\a_t\|_\infty\le1$, so the server sign is valid. Note that for MVSM-v2, we have  $\v_t=\operatorname{S_1}(\a_t)$ as well as
\begin{align}
 \E[\v_t]
 &
 =\E[\a_t]
 =\frac{\bar\v_t}{G}.\label{eq:composed-mean}
\end{align}
Consequently, the noise of the two-stage communication compression $ \e_t=\v_t-\bar\v_t/G$ satisfies
\begin{align}
 \E[\e_t]&=0,&
 \E[\|\e_t\|^2]
 &=d-\|\bar\v_t/G\|^2\le d.\label{eq:composed-variance}
\end{align}
Next, we introduce a virtual-iterate.
Write $c_\beta=\eta(1-\beta)/(G\beta)$ and, for $t\ge2$, define
\begin{align}
 \y_t=\x_t-c_\beta\bar\v_{t-1}.\label{eq:virtual-point}
\end{align}
The averaged momentum obeys $\bar\v_t=(1-\beta)\bar\v_{t-1}+\beta\bar\g_t$ for $t\ge2$. Since $\eta/G+c_\beta=\eta/(G\beta)$,
\begin{align}
 \y_{t+1}
 &=\x_t-\eta(\bar\v_t/G+\e_t)-c_\beta\bar\v_t\notag\\
 &=\x_t-\frac{\eta}{G\beta}
       \big((1-\beta)\bar\v_{t-1}+\beta\bar\g_t\big)-\eta\e_t\notag\\
 &=\y_t-\frac\eta G\bar\g_t-\eta\e_t.\label{eq:virtual-recursion}
\end{align}
Thus the virtual point takes an unbiased stochastic gradient step plus communication noise. 

Moreover, we ensure that
\begin{align}
 \|\y_t-\x_t\|\le c_\beta G\sqrt d
 =\eta\sqrt d\frac{1-\beta}{\beta}.\label{eq:virtual-displacement}
\end{align}
First, we consider the first step. Since $\bar\v_1=\bar\g_1$, the initialization must be treated separately:
\begin{align*}
 \y_2-\x_1=-\eta\v_1-c_\beta\bar\g_1.
\end{align*}
By equation~(\ref{eq:composed-mean}) and gradient unbiasedness,
\begin{align*}
 \E\langle\nabla f(\x_1),\y_2-\x_1\rangle
 =-\frac{\eta}{G\beta}\|\nabla f(\x_1)\|^2.
\end{align*}
Also $\|\v_1\|=\sqrt d$ and $\|\bar\g_1\|\le G\sqrt d$, so
\begin{align*}
 \|\y_2-\x_1\|\le\eta\sqrt d+c_\beta G\sqrt d
 =\frac{\eta\sqrt d}{\beta}.
\end{align*}
Smoothness therefore implies
\begin{align}
 \E f(\y_2)
 \le f(\x_1)-\frac{\eta}{G\beta}\|\nabla f(\x_1)\|^2
       +\frac{L\eta^2d}{2\beta^2}.\label{eq:v2-init}
\end{align}
Then, we analyze for every later step.
For $t\ge2$, smoothness  and equation~(\ref{eq:virtual-recursion}) give
\begin{align}
 \E[f(\y_{t+1})]
 \le f(\y_t)-\frac\eta G\langle\nabla f(\y_t),\nabla f(\x_t)\rangle
   +\frac{L\eta^2}{2}
       \E[\|\bar\g_t/G+\e_t\|^2].\label{eq:v2-descent-start}
\end{align}
Assumption~\ref{bg1} then yields
\begin{align}
 \E[\|\bar\g_t/G+\e_t\|^2]
 &=\E\left[\frac{\|\bar\g_t\|^2}{G^2}
          +d-\frac{\|\bar\v_t\|^2}{G^2}
          \right]\le2d.\label{eq:v2-second-moment}
\end{align}
Note that $-\langle\u,\w\rangle\le-\|\w\|^2/2+\|\u-\w\|^2/2$.
Applying this with $\u=\nabla f(\y_t)$, $\w=\nabla f(\x_t)$, and using smoothness and equation~(\ref{eq:virtual-displacement}), we obtain
\begin{align}
 \E f(\y_{t+1})
 \le\E f(\y_t)-\frac\eta{2G}\E\|\nabla f(\x_t)\|^2
    +L\eta^2d
    +\frac{L^2\eta^3d}{2G}\left(\frac{1-\beta}{\beta}\right)^2.
 \label{eq:v2-descent}
\end{align}
Since $\beta\le1$, the descent coefficient $\eta/(G\beta)$ in the first step is at least $\eta/(2G)$. Summing~\eqref{eq:v2-init} and~\eqref{eq:v2-descent} over all $T$ steps, and using $f(\y_{T+1})\ge f_*$, gives
\begin{align*}
 \frac\eta{2G}\sum_{t=1}^T\E\|\nabla f(\x_t)\|^2
 \le\Delta_f+\frac{L\eta^2d}{2\beta^2}
      +(T-1)L\eta^2d
      +\frac{(T-1)L^2\eta^3d}{2G}
         \left(\frac{1-\beta}{\beta}\right)^2.
\end{align*}
For $\eta=\gamma/\sqrt{dT}$,
\begin{align*}
 \frac1T\sum_{t=1}^T\E\|\nabla f(\x_t)\|^2
 \le\left(\frac{2G\Delta_f}{\gamma}+2LG\gamma\right)\sqrt{\frac dT}
       +\frac{L^2\gamma^2}{T}\left(\frac{1-\beta}{\beta}\right)^2
       +\frac{LG\gamma\sqrt d}{\beta^2T^{3/2}}
       \le\mathcal{O}\left(\sqrt{\frac dT}\right).
\end{align*}
The output index $\tau$ is independent and uniform on $\{1,\ldots,T\}$. Hence Jensen's inequality gives
\begin{align*}
 \E\|\nabla f(\x_\tau)\|_2
 \le\left(\frac1T\sum_{t=1}^T\E\|\nabla f(\x_t)\|^2\right)^{1/2}
 =\mathcal O(d^{1/4}T^{-1/4})
\end{align*}
for every fixed $\beta\in(0,1]$ and $\gamma>0$. 

\end{document}
